\section{Discussion}
\label{sec:discussion}

\subsection{The QK--OV Dissociation}
\label{sec:discussion_dissociation}

Our most striking finding is the dissociation between attention patterns (QK circuit) and output behavior (OV circuit) across data conditions.
On random data, the QK circuit produces strong induction attention patterns, and the OV circuit copies the attended token.
On natural text, the QK circuit distributes attention across many contextually relevant positions---weakening the induction-specific signal---but the OV circuit actually copies \emph{more effectively} when the attended content is natural language.

This dissociation suggests that the QK and OV circuits serve partially different functions depending on data distribution.
The QK circuit in induction heads may function as a general ``find relevant prior context'' mechanism, of which the $[A][B]\ldots[A] \rightarrow B$ pattern is a special case that dominates on random data but is one of many useful patterns on natural text.
Meanwhile, the OV circuit appears tuned to the statistics of natural language, producing stronger copying signals when the input matches the training distribution.

\subsection{Why Are Induction Scores So Much Lower on Natural Text?}
\label{sec:discussion_why}

Three factors likely contribute to the 10--40$\times$ reduction in induction scores on natural text:

\para{Attention competition.}
On random repeated sequences, the induction-pattern position is the single most informative attention target.
On natural text, many positions carry useful information---syntactic heads, semantic associates, positional neighbors---creating competition that dilutes the induction-specific signal.

\para{Multiple prediction strategies.}
On random data, copying is the only viable prediction strategy.
On natural text, the model can leverage syntactic patterns, semantic associations, n-gram statistics, and many other cues, reducing reliance on any single mechanism.

\para{Distributional mismatch.}
The random-sequence setup, by forcing exact token repetition in a fixed pattern, creates a maximally favorable environment for the induction mechanism.
Natural text, with its Zipfian token frequencies and rich contextual structure, is a far more complex prediction problem.

Notably, we ruled out one intuitive explanation: the difference is \emph{not} primarily driven by fewer repeated tokens in natural text.
Natural text has 41\% repeated tokens versus 50\% by construction in random sequences---a modest gap that cannot explain the 10--40$\times$ score difference.
We also ruled out matching precision: fuzzy (embedding-similarity) matching does not recover higher scores, confirming that the bottleneck is attention competition, not token-identity constraints.

\subsection{Implications for Mechanistic Interpretability}
\label{sec:discussion_implications}

Our results have practical implications for how the field studies induction heads and, more broadly, how we evaluate circuit-level mechanisms:

\para{Detection is reliable; quantification is not.}
Random-sequence detection correctly identifies \emph{which} heads perform induction (high rank correlation, high top-$K$ overlap).
However, the \emph{magnitude} of the induction score on random data dramatically overstates the strength of the induction attention pattern on natural text.
Researchers should complement random-sequence detection with naturalistic evaluation when making claims about the importance of specific heads.

\para{QK and OV circuits should be evaluated separately.}
The standard induction score conflates attention patterns and output behavior.
Our results show these can diverge: a head may show weak induction attention on natural text while having a strongly distribution-sensitive OV copying circuit.
Evaluating QK and OV contributions independently, as advocated by \citet{elhage2021mathematical}, provides a more complete picture.

\para{Distribution matters for circuit evaluation.}
More broadly, our findings caution against evaluating circuits exclusively on synthetic data.
The behavior of a circuit on artificial inputs may differ qualitatively from its behavior on the data distribution the model was trained on.

\subsection{Limitations}
\label{sec:limitations}

\para{Single model.}
We study only \gpttwo (124M parameters).
Larger models with more layers and heads may develop more specialized induction mechanisms, including genuinely separate naturalistic-only heads.
Scaling this analysis to GPT-2-medium, GPT-2-large, and the Pythia model suite is an important next step.

\para{Single detection metric.}
The induction score measures one specific attention pattern.
Other induction-like behaviors---attending to semantically similar contexts \cite{ren2024identifying}, computing conditional probabilities \cite{edelman2024evolution}, or performing abstract pattern completion---may require different detection methods that we do not explore here.

\para{Short sequences.}
Our 257-token sequences are relatively short.
Longer contexts may activate different attention patterns, particularly for late-layer heads that might benefit from richer context.

\para{Static analysis.}
We analyze a fully trained model at a single checkpoint.
The relationship between random and naturalistic induction behavior may evolve during training, with early-training heads potentially showing different cross-condition patterns than late-training heads \cite{olsson2022context, singh2024needs}.

\para{Threshold sensitivity.}
The 0.10 induction threshold is conventional but somewhat arbitrary.
Using percentile-based thresholds (top 10\% of heads) would yield slightly different categorizations, though our correlation-based analyses are threshold-independent.
